\documentclass{exam}
\usepackage[utf8]{inputenc}
\usepackage[T1]{fontenc}
\usepackage{geometry}
\usepackage{graphicx}
\usepackage{float}
\usepackage{hyperref}
\usepackage[english,indonesian]{babel}

\title{Optimasi Pemrosesan Template LaTeX Menggunakan Pendekatan Berbasis AI}
\author{Academic Support Agent}
\date{\today}

\begin{document}
\maketitle

\begin{abstract}

\end{abstract}

\chapter{PENDAHULUAN}


\chapter{TINJAUAN PUSTAKA}


% Original content (commented out):
% %
% This is a template for the MS Entrance Exam at National Tsing Hua University.  The university only provdes a Microsoft Word template, which is difficult to jointly collaborate and integrate multiple exam writers' questions. This LaTeX template attempts to approximate the official Word template but may require more fine tuning. It uses the exam documentclass, which can provide a rich collection of question formats.
%
%It can also be compiled with the answers option to show the solution key. Alternatively, you can also do \printanswers at the beginning of the exam.
%
\documentclass[a4paper,addpoints,12pt]{exam}
\usepackage[utf8]{inputenc}
\usepackage[english]{babel}
\usepackage[pdftex]{graphicx}
\usepackage{CJKutf8}
\usepackage{times}
\usepackage{mathptmx}
\usepackage{inconsolata}
\newcommand{\ctext}[1]{\begin{CJK}{UTF8}{bkai}#1\end{CJK}}
\pagestyle{head}
\extraheadheight{2.25cm}
\extrawidth{-1cm}
\header{\fontsize{14}{16}\selectfont~\\~\\~\\\ctext{系所班組別：資訊工程學系}\\
\ctext{考試科目（代碼）：2302計算機系統}
}{\ctext{\fontsize{16}{20}\selectfont 國立清華大學111學年度碩士班考試入學試題}\\~\\~\\}
              {~\\~\\~\\~\\~\\\vspace{12pt}\ctext{\fontsize{10}{12}\selectfont 共{\underline{\numpages}}頁，第{\underline{\thepage}}頁~~~  *請在【答案卷、卡】作答}}
\renewcommand{\choiceshook}{%
\setlength{\leftmargin}{1.85em}%
%\setlength{\labelwidth}{-\labelsep}%
}
\renewcommand{\choicelabel}{(\thechoice)}

\begin{document}

\begin{center}
\ctext{
\begin{tabular}{|c|c|c|c|c|}\hline
 題號    &  每題計分 & 選別 & 是否倒扣? & 備註\\\hline
  1 $\sim$ 40   &  2.5 分 & 單選 & 是，答錯扣率: 1/5 (即0.5分) & \\\hline
\end{tabular}}
\end{center}



\begin{questions}
%#1
\question[5]	Where is the main campus located? 

\begin{choices}
\CorrectChoice	City of Hsinchu (\ctext{新竹市}).
\choice	City of Taipei (\ctext{台北市}). 
\choice	New York City. 
\choice	Amsterdam. 
\choice	San Francisco.
\end{choices}

%#2
\question[5] Who was the President of NTHU in 2004? 

\begin{oneparchoices}
\choice	\ctext{陳力俊}.
\choice	\ctext{陳文村}. 
\choice	\ctext{賀陳弘}. 
\choice	\ctext{梅貽琦}. 
\CorrectChoice	\ctext{劉炯朗}. 
\end{oneparchoices}

%#3.
\question[10] Are electric scooters allowed on campus? Why or why not?

\vspace{2in}

\question[10] Compute $\displaystyle\int_0^1 x^2\,dx$

\makeemptybox{2in}

\question[10] Tell us about yourself.

\fillwithlines{1in}

\question[10] Tell us why you are applying to NTHU.

\fillwithdottedlines{1in}

\question[20] What if there were no trees at NTHU?
\begin{parts}
\part[5]	Would NTHU be indistinguishable from NYCU? 
\part[5]	Would all of NTHU be under water? 
\part[5]    Would NTHU be all grass? 
\part[5]    Would there be anybody left? 
\end{parts}

%%%%%%%%%%%
% This allows you to redefine the word for "points" if you want to show this in another language other than English by default.
\pointpoints{\ctext{分}}{\ctext{分}}

%#4
\question[5]	What is your favorite building? 
\begin{choices}
\choice	Delta Building. 
\CorrectChoice	Macronix Library. 
\choice	General Administration Building. 
\choice	Nuclear reactor. 
\choice	TSMC Building.
\end{choices}

%#5
\question[10] \fillin[National Tsing Hua University][2in] is the full name of NTHU.


\fullwidth{The following are True/False Questions. The fullwidth command can be given to provide instructions for the next set of questions without the indentation problem.}

\newcommand{\tf}[1][{}]{%
\fillin[#1][0.25in]%
}

\question[10] In the space below, plot the trajectory of a bouncing ball.

\fillwithgrid{1in}



\end{questions}

\end{document}


\end{document}